% Three birdcage capacitor topologies, each drawn as one ladder cell:
% two rungs (vertical) joined by top & bottom end-ring segments (horizontal).
%  high-pass: C in end-rings, L in rungs
%  low-pass : L in end-rings, C in rungs
%  band-pass: C and L in end-rings, C in rungs
\documentclass{article}
\usepackage[paperwidth=120cm,paperheight=120cm,margin=10mm]{geometry}
\pagestyle{empty}
\usepackage{fontspec}
\usepackage[bidi=basic]{babel}
\babelprovide[import]{english}
\babelprovide[main, import]{persian}
\babelfont{rm}[Renderer=Node, Path=../fonts/, Extension=.ttf, UprightFont=*-Regular, BoldFont=*-Bold]{Vazirmatn}
\babelfont{sf}[Renderer=Node, Path=../fonts/, Extension=.ttf, UprightFont=*-Regular, BoldFont=*-Bold]{Vazirmatn}
\providecommand{\setRTL}{}\providecommand{\setLTR}{}
\providecommand{\RL}[1]{\foreignlanguage{persian}{#1}}
\providecommand{\LR}[1]{\babelsublr{#1}}
\usepackage{tikz}
\usepackage{pgfplots}
\pgfplotsset{compat=1.18}
\usetikzlibrary{arrows.meta, positioning, shapes.geometric, shapes.misc, calc, fit, backgrounds, decorations.pathmorphing, patterns, angles, quotes}

% horizontal capacitor centered at (x,y), half-gap given; draws wires to ends
% #1 left x  #2 right x  #3 y  #4 label
\newcommand{\hcap}[4]{%
  \pgfmathsetmacro{\cx}{(#1+#2)/2}%
  \draw[thick] (#1,#3) -- (\cx-0.09,#3);
  \draw[thick] (\cx-0.09,#3+0.16) -- (\cx-0.09,#3-0.16);
  \draw[thick] (\cx+0.09,#3+0.16) -- (\cx+0.09,#3-0.16);
  \draw[thick] (\cx+0.09,#3) -- (#2,#3);
  \node[font=\footnotesize, above] at (\cx,#3+0.16) {#4};
}
% horizontal inductor between (x1,y) and (x2,y)
\newcommand{\hind}[4]{%
  \pgfmathsetmacro{\ia}{#1+0.12}%
  \pgfmathsetmacro{\ib}{#2-0.12}%
  \draw[thick] (#1,#3) -- (\ia,#3);
  \draw[thick,decorate,decoration={coil,aspect=0.5,segment length=3.6pt,amplitude=3.6pt}] (\ia,#3) -- (\ib,#3);
  \draw[thick] (\ib,#3) -- (#2,#3);
  \node[font=\footnotesize, above] at ({(#1+#2)/2},#3+0.16) {#4};
}
% vertical capacitor between (x,y1 top) and (x,y2 bottom)
\newcommand{\vcap}[4]{%
  \pgfmathsetmacro{\cy}{(#2+#3)/2}%
  \draw[thick] (#1,#2) -- (#1,\cy+0.09);
  \draw[thick] (#1-0.16,\cy+0.09) -- (#1+0.16,\cy+0.09);
  \draw[thick] (#1-0.16,\cy-0.09) -- (#1+0.16,\cy-0.09);
  \draw[thick] (#1,\cy-0.09) -- (#1,#3);
  \node[font=\footnotesize, right] at (#1+0.18,\cy) {#4};
}
% vertical inductor between (x,y1 top) and (x,y2 bottom)
\newcommand{\vind}[4]{%
  \pgfmathsetmacro{\ja}{#2-0.12}%
  \pgfmathsetmacro{\jb}{#3+0.12}%
  \draw[thick] (#1,#2) -- (#1,\ja);
  \draw[thick,decorate,decoration={coil,aspect=0.5,segment length=3.6pt,amplitude=3.6pt}] (#1,\ja) -- (#1,\jb);
  \draw[thick] (#1,\jb) -- (#1,#3);
  \node[font=\footnotesize, right] at (#1+0.18,{(#2+#3)/2}) {#4};
}

\begin{document}
\setRTL
\babelsublr{\begin{tikzpicture}[>=Stealth, font=\small]

  % geometry of one cell: left rung x=0.4, right rung x=2.4, top y=1.8, bot y=0
  % ============ HIGH-PASS ============
  \begin{scope}[shift={(0,0)}]
    \node[font=\bfseries] at (1.4,2.55) {بالاگذر};
    \node[font=\footnotesize] at (1.4,2.2) {(C در حلقه‌های انتهایی)};
    % top end-ring: C
    \hcap{0.4}{2.4}{1.8}{$C_e$}
    % bottom end-ring: C
    \hcap{0.4}{2.4}{0}{$C_e$}
    % rungs: L
    \vind{0.4}{1.8}{0}{$L_\ell$}
    \vind{2.4}{1.8}{0}{$L_\ell$}
  \end{scope}

  % ============ LOW-PASS ============
  \begin{scope}[shift={(4.4,0)}]
    \node[font=\bfseries] at (1.4,2.55) {پایین‌گذر};
    \node[font=\footnotesize] at (1.4,2.2) {(C در میله‌ها)};
    \hind{0.4}{2.4}{1.8}{$L_e$}
    \hind{0.4}{2.4}{0}{$L_e$}
    \vcap{0.4}{1.8}{0}{$C_\ell$}
    \vcap{2.4}{1.8}{0}{$C_\ell$}
  \end{scope}

  % ============ BAND-PASS ============
  \begin{scope}[shift={(8.8,0)}]
    \node[font=\bfseries] at (1.4,2.55) {میان‌گذر};
    \node[font=\footnotesize] at (1.4,2.2) {(C در هر دو)};
    % top end-ring: C then L
    \hcap{0.4}{1.25}{1.8}{$C_e$}
    \hind{1.25}{2.4}{1.8}{$L_e$}
    % bottom end-ring: C then L
    \hcap{0.4}{1.25}{0}{$C_e$}
    \hind{1.25}{2.4}{0}{$L_e$}
    % rungs: C
    \vcap{0.4}{1.8}{0}{$C_\ell$}
    \vcap{2.4}{1.8}{0}{$C_\ell$}
  \end{scope}

  % common note
  \node[font=\footnotesize, align=center] at (6.6,-0.65)
    {شاخه‌های عمودی $=$ میله‌ها \quad$|$\quad شاخه‌های افقی $=$ بخش‌های حلقهٔ انتهایی};

\end{tikzpicture}}
\end{document}
